%!TEX root = ../template.tex
%%%%%%%%%%%%%%%%%%%%%%%%%%%%%%%%%%%%%%%%%%%%%%%%%%%%%%%%%%%%%%%%%%%%
%% chapter1.tex
%% NOVA thesis document file
%%
%% Chapter with lots of dummy text
%%%%%%%%%%%%%%%%%%%%%%%%%%%%%%%%%%%%%%%%%%%%%%%%%%%%%%%%%%%%%%%%%%%%

\typeout{NT FILE chapter1.tex}%

\chapter{Introduction}
\label{cha:intro}

Defending static structures against fast flying objects is an age-old problem that appears in a variety of scenarios, such as in aerospace (i.e. space debris endangering orbiting stations) and in military. It is in this latter scenario that the problem is more pressing and more challenging as the objects, in this case referred to as missiles, are both self-propelled and guided. A particular area of interest is active defense, where an unmanned autonomous vehicle is driven to intercept and neutralize the missile. Many technologies have been developed to this end, however these missile defense platforms often rely on specialized machinery that is either difficult to deploy \textit{en-masse} or impractical in urban settings. Indeed, the twenty-first century has seen a rise in the need of missile defense in urban settings. At the same time there has also been a fast development of the capacities of commercial UAVs, which are inexpensive, lightweight, and easy to deploy when compared to traditional missile defense platforms. While much research has been done towards this end, resulting in a host of flight controllers that grant quad-rotor UAVs missile defense-capability, most methods do not provide early safety guarantees or trajectory feasibility analysis. A relatively novel method that may fill these gaps is reachability analysis. \\

Thus, the goal of this study is to design a flight controller for a quad-rotor UAV such that it defends a critical structure from a missile heading towards it, whether by intercepting it or deviating the missile's trajectory from the structure. The strategy considered for the controller design will be based on reachability tools, which utilises the concept of a reachable set. \\

\section{State of the Art}
\label{sec:state_of_the_art}

The topics relevant to this project include that of modelling \gls{UAVs}, specifically of missiles and quad-rotor drones; control techniques for guided missiles, specifically with stationary targets, and with adversary avoidance; control techniques for quad-rotor drones, specifically with non-cooperative intercept targets. Therefore it's necessary to map out the state of the art on these topics. \\

\subsection{Simulation}
\label{sec:simulation}

It would be impractical to iterate the design of such a control function based on real flight tests alone, specially in the early stages of the project, as each change would require a new, lengthy flight test. Therefore it is necessary to simulate the scenario under study, meaning, to simulate the missile and drone dynamics such as the aerodynamic effects acting on the vehicles, in response to input actuations. An almost ubiquitous tool used in aerodynamical simulation and control design are MATLAB/Simulink ~\cite{MATLAB}, thanks to the diversity of applications and efficiency they provide. Other tools such as X-Plane and Gazebo-ROS provide a visual simulation environment with MATLAB integration. Study ~\cite{MuhammadT} details a systematic methodology of running simulation techniques that sequentially integrate more and more of the actual system into the simulation. These are \gls{MILS}, \gls{SILS}, \gls{PILS} and \gls{HILS}, near-actual simulation techniques that can help speed up the flight system's development. 

\subsection{Modelling}
\label{sec:modelling}

\subsubsection{UAV Dynamics}

Obtaining an accurate mathematical model for the vehicle dynamics is crucial, whether it be for simulation or designing the flight systems. Regarding the modelling of missile dynamics, the most relevant phases of flight for this study are the midflight and homing phases. The rotational dynamics of the missile may be parametrized with Euler angles or quaternions.
When it comes to quad-rotor drone dynamics, this study is most concerned with the post-launch flight phase. In ~\cite{XiadongZ}, it is shown that the 6DoF dynamic model of a quad-rotor are complicated, strongly nonlinear, inherently unstable with coupled system variables and under-actuated. However, they may be thought of as two distinct control problems; the outer loop tracking the position and the inner loop tracking the attitude, which should simplify the controller design process. \\

\begin{figure}[htbp]
  \centering
  \includegraphics[width=0.8\linewidth]{ch1_UAV_diag}%
  \caption{Diagram of UAV model. Source:~\cite{XiadongZ}}
  \label{fig:UAV_diag}
\end{figure}

\subsubsection{Aerodynamic Effects}

Both the missile and the quad-rotor are expected to move at very high speeds. Therefore, it's clear that aerodynamic effects such as air drag must be considered in the modelling process. In the case of quad-rotors, blade flapping and thrust variations pose a significant disturbance to the system, and must also be accounted for. In ~\cite{HaomiaoH}, these effects are explicitly compensated for in the control loop of a quad-rotor drone.

\subsection{Control Systems}
\label{sec:control_systems}

It's worthwhile to review existing flight control systems, for both missiles and drones. The most widely used methodologies for missile control systems are \gls{PN} guidance, and for UAVs they are Optimal Control-based. These may further be classified into Linear Control Systems and Model Based Nonlinear Control Systems ~\cite{FaridK}. \\

\subsubsection{GNC}

The systems that control the movements of UAVs can be described by their three core functions: \gls{GNC}. Guidance refers to the systems that determine the desired paths from the vehicle's current position to the target waypoint, as well as establishing the changes in velocity, rotation or acceleration to achieve this course. Navigation refers to the systems that determine, at a given instant of time, the vehicle's location, velocity, rotation, and angular rate. These variables are the system's state vector. Control is referring to the manipulation of forces and torques, by means of actuators (i.e., steering device, thrusters, etc.), needed to follow the desired path while maintaining vehicle stability ~\cite{PesceV}. The goal of this project is, essentially, to design a control function that given the current state vector of both the missile and the quad-rotor drone, will drive the drone to intercept the missile. \\

\begin{figure}[htbp]
  \centering
  \includegraphics[width=0.5\linewidth]{ch1_GNC}%
  \caption{The overall structure of guidance, navigation, and control systems onboard a RUAS. Source:~\cite{FaridK}, Figure 3}
  \label{fig:GNC}
\end{figure}

\subsubsection{Linear Control Methods}

Linear control methods are the least computationally complex methods and the easiest to implement, and generally consist in driving the controlled system variable to converge to a desired setpoint, based on negative feedback. However, they tend to have limited robustness when the desired value is variable, or when the vehicle is under-acted, this is, when there are fewer actuators than the number of degrees of freedom to be controlled. To overcome these limitations, the controlled system must be linearized, usually by finding linear dynamics around a stable operating point or trajectory.

\paragraph*{PID}
One of the most common Linear control methods is the \gls{PID} control. PID is based on negative feedback, and it aims to converge the controlled variable's value to a given setpoint, according to given specifications. From these specifications the proportional, integral, and derivative gain are selected. These will be multiplied by the respective parts of the error signal, meaning the difference between the desired setpoint and the current value of the controlled variable. This transformed signal is injected back into the system as input. In~\cite{SanchezI} a survey of the performance of different PID methods on UAVs is conducted. It reveals that linear PIDs are effective at hovering, moving to a waypoint and following a trajectory, with the advantages of being simple to design and implement as well as not needing knowledge of the system model. It has several disadvantages however, such as requiring the model to be linearized, it is not robust to external disturbances and it's only valid only a specific operation point under specific conditions.

\paragraph*{LQR}
Another commonly used linear control method is the \gls{LQR}. The system dynamics are described by a set of linear differential equations and the cost is described by a quadratic function, usually defined as the geometric norm of the difference between the state vector and the setpoint at the current instant. Weighting factors are selected then applied to different parts of the cost function, such as the final value or the sum of all deviations. The control actuation may also feature in the cost function. The LQR provides the optimal control law given these parameters and is robust to perturbations. Moreover, it can handle time-varying models. In ~\cite{ShahidaK}, a quad-rotor drone is controlled with an LQR based controller. It is shown that this controller is robust, with low steady state error.

\paragraph*{Gain Scheduling}
Gain Scheduling is another linear control method, in which non-linear system dynamics are modelled as a collection of simplified linear models, with each model representing a particular operating regime. Each model has a corresponding linear control law. In ~\cite{JingQ}, a gain scheduling-based controller is implemented in a quad-rotor drone, where different PID controllers are used depending on the current operating regime. The study shows that, relative to the singular PID controller, the performance in trajectory tracking is increased, while maintaining the advantages of PIDs. \\

\subsubsection{Nonlinear Control Methods}

Model-Based Nonlinear Controllers are another family of flight control systems. They have the advantage of being able to handle systems'  nonlinear dynamics, making them more robust and versatile than linear control methods. This comes at the cost of a more complex, more computationally expensive implementation.

\paragraph*{MPC}
\gls{MPC} as a model-based control method is an extension of the LQR. However, the MPC estimates the model output in a finite time horizon then minimizes tracking error for this horizon. Usually, the control law is updated online. Designing the MPC controller requires an explicit linear model of the system, and this is not usually the case in UAV control, making linearization necessary. Study ~\cite{AtharvaG} considers both a UAV dynamic target interception scenario and a surveillance scenario, and proposes a novel MPC design with variable time steps. The controller was tested both in simulation and real-world flight tests, in 2D and 3D space, and was shown to be effective and efficient. 

\paragraph*{NMPC}
\gls{NMPC} is a variation of the linear MPC that can handle nonlinear models. The optimization problem is usually non-convex, therefore this method does not usually yield a closed-form expression and requires numerical methods to obtain the optimal control law for this time step. This introduces the "curse of dimensionality", where the complexity grows exponentially with the horizon. Study ~\cite{DoraN} tackles the problem of collision avoidance in a multi-UAV agricultural scenario. It compares the performance of three NMPC-based control strategies: distributed NMPC where the NMPCs of the UAVs are coupled through nonlinear constraints, distributed NMPC where the NMPCs of the UAVs are coupled through penalty costs, and imposing flight corridors by modelling uncertainties and disturbances. It shows that in a scenario with no path intersections the flight corridor strategy ranks better, however it fails to ensure collision avoidance when there are path intersections, while the other two methods are generally effective but with varying performances.

\paragraph*{Backstepping}
Backstepping Control is a nonlinear control method, commonly used in systems with a relative degree higher than one. This method separates the system into various subsystems such that the relative degree of each is one, then a "virtual control law" is applied to the higher-order subsystem such that it stabilizes the lower order one. This creates a series of "virtual controllers" that ultimately allow the controlled variables to converge towards the setpoint. This allows for an iterative, systematic design process that especially favors the layered structure of UAV dynamics. However, the complexity of these controllers rises dramatically as the relative degree increases, since each differentiation step is increasingly complicated to solve analytically. Study \cite{QingboG} compares the performances of two route-planning approaches for a UAV controller with collision avoidance based on backstepping: a safety ball approach where the UAV stays a minimum distance away from obstacle boundaries while converging towards the target, and a point-mass approach that treats obstacles as a sphere with radius equal to the obstacle's largest dimension. It shows that the safety ball approach has better performance when it comes to computational complexity, but does not minimise the tracking error.

\paragraph*{CLF-CBF}
\gls{CLF}, along with \gls{CBF} and unified by \gls{QP} establish their own class of nonlinear controllers. The CLF uses the Lyapunov function, which represents an objective. With the CLF, the input norm is minimized while subject to an inequality constraint on the derivative of a Lyapunov function, and the system is guaranteed to converge exponentially towards the objective. CLF methodology allows the system to have several objectives simply by adding CLFs representing them to the constraints. CBFs are similar to CLFs, except they use Lyapunov-like functions called Barrier functions to define a boundary that the system cannot cross, and will diverge away from ~\cite{AmesA}. By joining CLFs and CBFs, a powerful, a robust controller is achieved. Other methods may be used along with CLF/CBF controllers to further expand their efficacy, such as backstepping. The study ~\cite{LiB} proposes a defense-winning strategy for active target protection involving multiple nonlinear players, utilising a CBF-QP controller. \\

\begin{figure}[htbp]
  \centering
  \includegraphics[width=0.3\linewidth]{ch1_CBF}%
  \caption{ Illustration of the m-sphere enclosure region. (a) 2-sphere enclosure region in a 2D plane. (b) 3-sphere enclosure region in the 3D space. Source:~\cite{LiB}, Figure 1}
  \label{fig:CBF}
\end{figure}

\subsubsection{Reachability-Based Guidance}
PN and optimal control based methods, while traditionally used have their limitations. PN guidance cannot account for changes in the adversary's behavior, while optimal control cannot guarantee feasibility given the current conditions. A modern alternative is Reachability based guidance, that aims to grapple with these limitations. Reachability analysis can be used to implicitly guarantee feasibility of a trajectory, or for collision avoidance. Study ~\cite{AkalinB} employs a novel method for determining the reachable set boundaries for aerodynamic interceptors.

\paragraph*{Approximate Geometric Method}
One way to calculate the reachable set is the Approximate Geometric Method. It approximates the reachable set of the model by manipulating convex shapes such as ellipsoids, zonotopes, or polytopes such that they closely fit the bounds of the expected set in the next time steps. This method can even be used with nonlinear models, as shown in ~\cite{SilvestreD1}. A drawback of these methods is that, as the set computation is iterated for each following time step, the dimensional complexity grows exponentially, thus increasing the computational burden of this method. However, study ~\cite{RegoF} proposes an algorithm that significantly reduces this effect. 

\begin{figure}[htbp]
  \centering
  \includegraphics[width=0.8\linewidth]{ch1_reachability}%
  \caption{Example attainable zone. Source:~\cite{RegoF}, Figure 2}
  \label{fig:reachability}
\end{figure}

